\chapter{Funções de uma variável real}\label{ch:b2:realfun}

Antes que a análise passe às funções de funções (\cref{ch:b2:funcseq}),
compensa conhecer a paisagem de uma variável com mais detalhe do que
o primeiro ano exigia: quão descontínua pode ser uma função monótona, quão
regular deve ser uma função convexa, e de que propriedades especiais gozam
as derivadas (Darboux). Esses resultados estruturais são curtos,
afiados e muito queridos pelos examinadores.

\section{Funções monótonas}

\begin{theorem}[Regularidade das funções monótonas]\label{thm:b2:realfun:monotone}
Seja $f \colon I \to \R$ crescente num intervalo.
\begin{enumerate}
  \item Em todo ponto interior $a$, existem os limites laterais:
        \[
        f(a^-) = \sup_{x < a} f(x) \;\leq\; f(a) \;\leq\;
        f(a^+) = \inf_{x > a} f(x) ;
        \]
        toda descontinuidade é um \emph{salto}.
  \item O conjunto das descontinuidades de $f$ é no máximo
        \omterm{def:b2:structures:countable}{enumerável}.\index{função monótona!descontinuidades}
\end{enumerate}
\end{theorem}

\begin{proof}
(1) O conjunto $\{f(x) : x < a\}$ é não vazio e majorado por $f(a)$:
seu supremo $s$ satisfaz $f(x) \to s$ quando $x \to a^-$ (dado
$\varepsilon$, algum $f(x_0) > s - \varepsilon$, e a monotonicidade
prende $f(x) \in \intoc{s - \varepsilon}{s}$ para $x \in
\intoo{x_0}{a}$). Simetricamente à direita.

(2) A cada descontinuidade $a$ associe o intervalo \omterm{def:b2:metric:topology}{aberto} não vazio
$J_a = \intoo{f(a^-)}{f(a^+)}$ (um salto genuíno). Para $a < b$
descontinuidades, $J_a$ e $J_b$ são disjuntos: $f(a^+) \leq f(c)
\leq f(b^-)$ para qualquer $c$ entre elas. Cada $J_a$ contém um racional;
descontinuidades distintas recebem racionais distintos: uma injeção do
conjunto das descontinuidades em $\Q$, que é \omterm{def:b2:structures:countable}{enumerável}
(\cref{prop:b2:structures:countablestable}).
\end{proof}

\begin{example}
A cota é ótima: fixe uma enumeração $(r_n)$ de $\Q \cap
\intoo{0}{1}$ e ponha $f(x) = \sum_{n : r_n \leq x} 2^{-n}$
(uma definição por \omterm{def:b2:series:summable}{família somável},
\cref{def:b2:series:summable}). Então $f$ é crescente em
$\intcc{0}{1}$ e descontínua exatamente em cada racional de
$\intoo{0}{1}$ (salto $2^{-n}$ em $r_n$): uma função monótona
\emph{pode} ser descontínua num \omterm{def:b2:structures:countable}{conjunto enumerável} denso.
\end{example}

\begin{example}[Os saltos não podem superar a subida]\label{ex:b2:realfun:jumpbudget}
Para $f$ crescente em $\intcc{a}{b}$, os saltos têm um orçamento:
se $a < c_1 < \dots < c_m < b$ são descontinuidades com saltos
$s_i = f(c_i^+) - f(c_i^-) > 0$, então, escolhendo pontos
intercalados $a < c_1 < t_1 < c_2 < \dots$ e usando a monotonicidade em
cada pedaço,
\[
\sum_{i=1}^{m} s_i \;\leq\; f(b) - f(a) :
\]
a subida total majora o salto total. Consequência: para
cada $k$, no máximo $k\,\bigl(f(b) - f(a)\bigr)$ descontinuidades
têm salto $\geq \frac1k$ --- um refinamento quantitativo do
\cref{thm:b2:realfun:monotone}~(2), pois o conjunto das descontinuidades
é a união \omterm{def:b2:structures:countable}{enumerável} sobre $k$ desses conjuntos finitos. Na
função de saltos racionais acima, o orçamento é gasto exatamente:
os saltos $2^{-n}$ somam $1 = f(1^+) - f(0^-)$ no sentido estendido
óbvio. As funções monótonas podem saltar densamente, mas apenas
com uma mesada estrita.
\end{example}

\section{Funções convexas}

\begin{lemma}[Desigualdade das inclinações]\label{lem:b2:realfun:slopes}
Seja $f$ convexa em $I$ e $x < y < z$ em $I$. Então
\[
\frac{f(y) - f(x)}{y - x}
\;\leq\; \frac{f(z) - f(x)}{z - x}
\;\leq\; \frac{f(z) - f(y)}{z - y} :
\]
as inclinações das cordas crescem nas duas extremidades.\index{função
convexa!desigualdade das inclinações}
\end{lemma}

\begin{proof}
Escreva $y = \frac{z - y}{z - x}\,x + \frac{y - x}{z - x}\,z$: uma
combinação convexa, pois os dois coeficientes são positivos e
somam $1$. A convexidade dá
\[
f(y) \;\leq\; \frac{z-y}{z-x}\,f(x) + \frac{y-x}{z-x}\,f(z).
\]
Para a desigualdade da esquerda, subtraia $f(x)$ dos dois lados, usando
$\frac{z-y}{z-x} - 1 = -\frac{y-x}{z-x}$:
\[
f(y) - f(x) \leq \frac{y - x}{z - x}\bigl(f(z) - f(x)\bigr),
\]
e divida por $y - x > 0$. Para a desigualdade da direita, subtraia
em vez disso de $f(z)$:
\[
f(z) - f(y) \geq f(z) - \frac{z-y}{z-x}f(x) -
\frac{y-x}{z-x}f(z)
= \frac{z - y}{z - x}\bigl(f(z) - f(x)\bigr),
\]
e divida por $z - y > 0$. Os dois passos exibidos são a mesma
identidade baricêntrica lida contra uma extremidade diferente.
\end{proof}

\begin{theorem}[Regularidade das funções convexas]\label{thm:b2:realfun:convexreg}
Seja $f$ convexa num intervalo $I$.
\begin{enumerate}
  \item Em todo ponto interior, $f$ tem derivadas laterais
        finitas $f'_g \leq f'_d$; ambas são funções crescentes
        do ponto; em particular, $f$ é \omterm{def:b2:metric:continuity}{contínua} no
        interior de $I$ (mas possivelmente não nas extremidades).
  \item $f$ fica acima de cada \emph{reta de suporte}: para $a$ interior
        e qualquer $m \in \intcc{f'_g(a)}{f'_d(a)}$,
        \[
        f(x) \geq f(a) + m(x - a) \qquad (x \in I).
        \]
  \item (Jensen, com pesos) Para $x_i \in I$ e pesos $\lambda_i
        \geq 0$, $\sum\lambda_i = 1$:
        \[
        f\Bigl(\sum_i \lambda_i x_i\Bigr) \leq \sum_i \lambda_i
        f(x_i) .\index{desigualdade de Jensen}
        \]
\end{enumerate}
\end{theorem}

\begin{proof}
(1) Fixe $a$ interior. Pelo~\cref{lem:b2:realfun:slopes}, a inclinação
$\tau(h) = \frac{f(a + h) - f(a)}{h}$ é uma função crescente de
$h$ (dos dois lados, e $\tau(h_-) \leq \tau(h_+)$ para $h_- < 0 <
h_+$). Logo $\tau$ tem limite finito quando $h \to 0^-$ (crescente,
majorada por qualquer inclinação à direita) --- esse é $f'_g(a)$ --- e quando
$h \to 0^+$ ($f'_d(a)$), com $f'_g(a) \leq f'_d(a)$. Derivadas laterais
finitas forçam a \omterm{def:b2:metric:continuity}{continuidade} em $a$. Monotonicidade no
ponto: para $a < b$ interiores, $f'_d(a) \leq \frac{f(b) - f(a)}{b -
a} \leq f'_g(b)$, de novo pela desigualdade das inclinações.

(2) Para $x > a$: $\frac{f(x) - f(a)}{x - a} \geq f'_d(a) \geq m$;
para $x < a$: $\frac{f(a) - f(x)}{a - x} \leq f'_g(a) \leq m$. As duas
se rearranjam na afirmação.

(3) Indução no número de pontos exatamente como no volume do primeiro
ano de graduação (o caso de dois pontos é a definição) --- ou de um golpe só:
aplique (2) em $a = \sum\lambda_i x_i$ e faça a média das desigualdades das retas
de suporte nos pontos $x_i$ com pesos $\lambda_i$: $\sum_i
\lambda_i f(x_i) \geq f(a) + m\sum_i\lambda_i(x_i - a) = f(a)$.
\end{proof}

\begin{omfigure}
\begin{tikzpicture}
\begin{axis}[omaxis, width=8.2cm, height=5.4cm,
    xmin=-2.4, xmax=2.5, ymin=-0.9, ymax=4.5,
    xtick={-1.5, 0.5, 2}, ytick=\empty]
  \addplot[omDef, very thick, domain=-2.1:2.15, samples=80]
    {x^2};
  \addplot[omThm, thick, domain=-1.5:2, samples=2]
    {0.5*x + 3};
  \addplot[omProp, thick, domain=-1.1:2.2, samples=2]
    {x - 0.25};
  \node[font=\small, omThm] at (axis cs:-1.05,3.15)
    {corda};
  \node[font=\small, omProp] at (axis cs:1.75,0.95)
    {reta de suporte};
  \node[font=\small, omDef] at (axis cs:-1.95,2.6) {$f$};
\end{axis}
\end{tikzpicture}

{\small A convexidade numa só figura: entre $-1.5$ e $2$ o
gráfico de $f(x) = x^2$ fica abaixo de sua corda (a definição) e
acima da reta de suporte em $x = 0.5$
(\cref{thm:b2:realfun:convexreg}~(2)) --- toda desigualdade do
problema de fim de semana deste capítulo é um rearranjo dessas duas
posições.}
\end{omfigure}

\begin{example}[Descontinuidade na extremidade]
Em $\intcc{0}{1}$, a função $f(0) = 1$, $f(x) = 0$ para $x > 0$,
é convexa mas descontínua na extremidade $0$: a afirmação (1) é
ótima.
\end{example}

\begin{example}[Bicos e o feixe de retas de suporte]\label{ex:b2:realfun:corner}
Para $f(x) = \abs x$ em $a = 0$: as derivadas laterais são
$f'_g(0) = -1$ e $f'_d(0) = +1$, e o
\cref{thm:b2:realfun:convexreg}~(2) entrega uma reta de suporte para
\emph{cada} inclinação $m \in \intcc{-1}{1}$:
\[
\abs x \geq m\,x \qquad (x \in \R,\ -1 \leq m \leq 1),
\]
cada uma sendo igualdade exatamente numa semirreta ou em $0$. Uma função
convexa é derivável em $a$ precisamente quando o feixe
colapsa numa única reta ($f'_g(a) = f'_d(a)$); os bicos carregam
um intervalo de tangentes. Esse feixe é o germe em dimensão finita
do \emph{subdiferencial} da otimização convexa ---
e a razão pela qual as funções convexas são tão robustas: mesmo onde a
derivada falha, a geometria de suporte sobrevive, e é só isso
que a demonstração de Jensen usou.
\end{example}

\begin{example}[Desigualdade das médias de potência]\label{ex:b2:realfun:powermeans}
Para $0 < p < q$ e $x_i$ positivos com pesos $\lambda_i$
somando $1$, aplicando Jensen à função convexa $t \mapsto t^{q/p}$
nos pontos $x_i^p$:
\[
\Bigl(\sum \lambda_i x_i^{p}\Bigr)^{1/p}
\leq \Bigl(\sum \lambda_i x_i^{q}\Bigr)^{1/q} :
\]
as médias de potência crescem com o expoente --- contendo MA--MQ e,
no limite $p \to 0$ (\cref{exo:b2:realfun:6}), a desigualdade
MA--MG mais uma vez.
\end{example}

\begin{omfigure}
\begin{tikzpicture}
\begin{axis}[omaxis, width=8.6cm, height=5.6cm,
    xmin=-8.5, xmax=8.5, ymin=0.5, ymax=4.5,
    xtick={-4,-1,1,2,4}, ytick={1,2,4},
    xlabel={$p$}, ylabel={$M_p$}]
  \addplot[omDef, very thick, domain=-8:-0.4, samples=90]
    {((1/3)*(1^x + 2^x + 4^x))^(1/x)};
  \addplot[omDef, very thick, domain=0.4:8, samples=90]
    {((1/3)*(1^x + 2^x + 4^x))^(1/x)};
  \addplot[omThm, dashed, domain=-8.5:8.5, samples=2] {4};
  \addplot[omThm, dashed, domain=-8.5:8.5, samples=2] {1};
  \node[font=\small, omDef] at (axis cs:-5.4,1.55)
    {$M_p(1,2,4)$};
  \node[font=\small, omThm] at (axis cs:-6.4,3.75)
    {$\max = 4$};
  \node[font=\small, omThm] at (axis cs:6.3,1.25)
    {$\min = 1$};
  \node[font=\small, omProp] at (axis cs:0.2,1.83) {$2$};
\end{axis}
\end{tikzpicture}

{\small A média de potência $M_p$ dos valores $1, 2, 4$ (pesos
iguais), como função do expoente $p$: crescente de
$\min = 1$ (quando $p \to -\infty$) até $\max = 4$ (quando $p \to
+\infty$), passando pelas médias harmônica ($p = -1$), geométrica (o vão
em $p = 0$, valor $2$), aritmética ($p = 1$) e quadrática ($p
= 2$). Toda a cadeia clássica das desigualdades entre médias é
uma única curva crescente --- demonstrada no problema de fim de semana
deste capítulo, Parte III.}
\end{omfigure}

\begin{example}[Entropia máxima]\label{ex:b2:realfun:entropy}
Para um vetor de probabilidade $(p_1, \dots, p_n)$ (positivo,
somando $1$), a entropia $H(p) = -\sum_i p_i\ln p_i$
satisfaz
\[
H(p) \leq \ln n ,
\qquad\text{com igualdade se e somente se } p_i = \frac1n \text{ para todo }
i .
\]
Demonstração por Jensen (\cref{thm:b2:realfun:convexreg}~(3)) aplicado à
função \emph{côncava} $\ln$ com pesos $p_i$ nos pontos
$\frac{1}{p_i}$:
\[
H(p) = \sum_i p_i\ln\frac{1}{p_i}
\leq \ln\Bigl(\sum_i p_i\,\frac1{p_i}\Bigr) = \ln n ,
\]
forçando a igualdade que todos os pontos $\frac1{p_i}$ sejam iguais (concavidade
estrita), isto é, $p$ uniforme. Equivalentemente, isso é
\cref{exo:b2:realfun:7} com $q$ uniforme. A incerteza é
maximizada pela ignorância uniformemente espalhada --- o princípio
variacional por trás da codificação, da mecânica estatística e das aparições
da entropia no~\cref{ch:b2:randomvar}.
\end{example}

\begin{method}[Achar a função convexa por trás de uma desigualdade]\label{met:b2:realfun:spotting}
A maioria das desigualdades clássicas é Jensen fantasiado; para
despi-lo: (1) normalize de modo que apareça uma \emph{média ponderada}
(pesos positivos, somando $1$ --- divida por uma massa total se
preciso); (2) veja que função é aplicada dentro e fora
da média: a afirmação ``$f(\text{média}) \leq$
média de $f$'' nomeia a convexa $f$; (3) certifique a convexidade pela
segunda derivada e trate a igualdade pela estrita; (4)
se nenhuma média está visível, tome logaritmos primeiro --- produtos
e potências viram médias, e a concavidade de $\ln$ carrega
MA--MG, Young e seus parentes (o problema de fim de semana deste capítulo
roda os passos 1--4 em cada um deles). Se nem os logaritmos revelam
uma média, tente ler a desigualdade como monotonicidade
das inclinações (\cref{lem:b2:realfun:slopes}) ---
enunciados de superaditividade como \cref{exo:b2:realfun:9} moram
ali.
\end{method}

\begin{remark}[Armadilhas comuns]\label{rem:b2:realfun:pitfalls}
(i) A convexidade não é preservada por produtos: $x$ e $(x - 1)^2$
são convexas em $\intcc{0}{2}$, mas seu produto $x(x-1)^2$ tem
segunda derivada $6x - 4$, negativa em
$\intco{0}{\frac23}$ --- não é convexo; nem a convexidade é preservada
por composição sem monotonicidade
(\cref{exo:b2:realfun:10}). (ii) Jensen se inverte para funções
côncavas: metade das desigualdades clássicas é a versão
côncava com $\ln$; aplicar a forma convexa a $\ln$ é o
modo mais rápido de demonstrar MA--MG \emph{de trás para a frente}. (iii) A
convexidade no ponto médio sozinha não implica convexidade --- é preciso \omterm{def:b2:metric:continuity}{continuidade} (ou
mera limitação) (\cref{exo:b2:realfun:8}); os
contraexemplos patológicos vivem além dos axiomas deste livro. (iv)
Uma função convexa num intervalo \emph{\omterm{def:b2:metric:topology}{aberto}} é \omterm{def:b2:metric:continuity}{contínua},
e até localmente \omterm{def:b2:metric:continuity}{lipschitziana} (\cref{exo:b2:realfun:12}); nas
extremidades, nada é de graça. (v) As derivadas obedecem a Darboux mas
não precisam ser \omterm{def:b2:metric:continuity}{contínuas} (\cref{ex:b2:realfun:oscillation}):
``$f'$ não tem saltos'' nunca significa ``$f'$ é \omterm{def:b2:metric:continuity}{contínua}''.
\end{remark}

\section{A propriedade de Darboux}

\begin{theorem}[Darboux]\label{thm:b2:realfun:darboux}
Seja $f$ derivável num intervalo $I$. Então $f'$ assume todo
valor entre dois quaisquer de seus valores --- ainda que $f'$ não precise ser
\omterm{def:b2:metric:continuity}{contínua}.\index{teorema de Darboux}
\end{theorem}

\begin{proof}
Sejam $a < b$ em $I$ e $v$ estritamente entre $f'(a)$ e $f'(b)$,
digamos $f'(a) < v < f'(b)$. A função $g(x) = f(x) - vx$ é
derivável com $g'(a) < 0 < g'(b)$: seu mínimo em
$\intcc{a}{b}$ (atingido: \omterm{def:b2:metric:continuity}{continuidade} num \omterm{def:b2:metric:compact}{compacto}) não está em $a$
(logo após $a$, $g$ decresce abaixo de $g(a)$) nem em $b$ (logo
antes de $b$, $g$ está abaixo de $g(b)$): ele é interior, e aí $g'(c)
= 0$, isto é, $f'(c) = v$. (Isso era um exercício com estrela do primeiro ano; seu
lugar na teoria é aqui.)
\end{proof}

\begin{example}[Que funções são derivadas?]\label{ex:b2:realfun:whichderivatives}
O teorema de Darboux é uma máquina de não existência. A função piso
$\lfloor x\rfloor$ não é a derivada de nenhuma função em
$\R$: ela assume os valores $0$ e $1$ mas pula $\frac12$ em
$\intcc{0}{1}$, o que o~\cref{thm:b2:realfun:darboux} proíbe para
derivadas. O mesmo veredicto atinge toda função com salto
--- sinal, Heaviside, todas as funções escada --- por mais inocentes
que pareçam; suas ``primitivas'' ($\abs x$ para o sinal, etc.)
existem apenas longe do salto e se atam ali com um bico.
Contraste: a função $f'$ selvagemente descontínua do
\cref{ex:b2:realfun:oscillation} \emph{é} uma derivada --- sua
descontinuidade é uma oscilação, que Darboux tolera. A
fronteira entre os dois comportamentos é exatamente o
corolário sem saltos a seguir.
\end{example}

\begin{corollary}\label{cor:b2:realfun:nojumps}
Uma derivada não tem descontinuidades de salto: se $f'(a^-)$ e
$f'(a^+)$ existem, eles valem $f'(a)$. As descontinuidades de uma
derivada são sempre do tipo oscilação (a derivada de $x^2\sin\frac1x$
em $0$, volume do primeiro ano de graduação).
\end{corollary}

\begin{proof}
Se $f'(a^+) = \lim_{x\to a^+} f'(x)$ existe e difere de
$f'(a)$, os valores estritamente entre eles seriam pulados por $f'$ numa
vizinhança à direita --- contradizendo Darboux nos intervalos
$\intcc{a}{a + h}$. (Alternativamente: o teorema do valor médio força
$f'(a) = \lim_{h\to0^+} \frac{f(a+h)-f(a)}{h} = f'(a^+)$, sendo o
quociente de diferenças um valor de $f'$ num ponto intermediário.)
O mesmo à esquerda.
\end{proof}

\begin{example}[A derivada oscilante canônica]\label{ex:b2:realfun:oscillation}
Seja $f(x) = x^2\sin\frac1x$ para $x \neq 0$ e $f(0) = 0$. Em
$0$: $\bigl|\frac{f(h) - f(0)}{h}\bigr| = \abs{h\sin\frac1h}
\leq \abs h \to 0$, logo $f'(0) = 0$ existe. Fora de $0$,
\[
f'(x) = 2x\sin\frac1x - \cos\frac1x ,
\]
cujo primeiro termo tende a $0$ enquanto $\cos\frac1x$ oscila
por $\intcc{-1}{1}$ em todo intervalo $\intoo{0}{\delta}$: o
limite $f'(0^+)$ não existe. Assim $f'$ está definida em toda parte mas é
descontínua em $0$ --- e, exatamente como
o~\cref{cor:b2:realfun:nojumps} prevê, a descontinuidade é uma
oscilação, e não um salto: em cada $\intoo{0}{\delta}$, $f'$ ainda
varre um intervalo inteiro em torno de $0$. As derivadas podem ser selvagens, mas
só do jeito compatível com Darboux.
\end{example}

\begin{remark}[Onde este capítulo é usado]
A convexidade é o motor da indústria de desigualdades: o
problema de fim de semana deste capítulo fabrica dela Young, Hölder, Minkowski
e a cadeia das médias de potência, que a teoria das \omterm{def:b2:nvs:norm}{normas} do
\cref{ch:b2:nvs} e as estimativas integrais do~\cref{ch:b2:integration} consomem;
Jensen reaparece em probabilidade como as desigualdades de momentos do
\cref{ch:b2:randomvar}. A regularidade monótona volta no
\cref{ch:b2:integration} (funções monótonas são integráveis) e,
no volume do terceiro ano de graduação, como a derivabilidade em quase toda parte
das funções monótonas --- onde ``uma quantidade \omterm{def:b2:structures:countable}{enumerável} de saltos'' se torna
o primeiro passo da teoria de Lebesgue.
\end{remark}

\section{Exercícios}

\begin{exercise}[$\star$]\label{exo:b2:realfun:1}
Determine os conjuntos de descontinuidade e os tamanhos dos saltos: $\lfloor x
\rfloor$; $\;x - \lfloor x\rfloor$; $\;\lfloor x \rfloor + \sqrt{x
- \lfloor x\rfloor}$; a função do exemplo que segue o
\cref{thm:b2:realfun:monotone} restrita aos racionais diádicos
$r_n$.
\end{exercise}

\begin{exercise}[$\star$]\label{exo:b2:realfun:2}
Prove que uma função crescente $f \colon I \to \R$ com a
propriedade do valor intermediário (a imagem de todo subintervalo é um
intervalo) é \omterm{def:b2:metric:continuity}{contínua}.
\end{exercise}

\begin{exercise}[$\star$]\label{exo:b2:realfun:3}
Quais das seguintes são convexas em seu domínio? $x \mapsto
x\ln x$ ($x > 0$); $\;x \mapsto \ln(1 + \eu^x)$; $\;x \mapsto
\sqrt{1 + x^2}$; $\;x \mapsto x^3$.
\end{exercise}

\begin{exercise}[$\star\star$]\label{exo:b2:realfun:4}
Seja $f$ convexa em $\R$ e majorada. Prove que $f$ é
constante. \emph{(Se $f(a) \neq f(b)$, a desigualdade das inclinações
propaga a inclinação não nula da corda: além do ponto de maior
valor, $f$ cresce ao menos linearmente --- contradizendo a limitação.
Trate os dois sinais da inclinação.)} Deduza que uma função convexa em
$\R$ com assíntota nas duas pontas é afim.
\end{exercise}

\begin{exercise}[$\star\star$]\label{exo:b2:realfun:5}
Seja $f$ derivável em $I$ com $f'$ monótona. Prove que
$f'$ é \omterm{def:b2:metric:continuity}{contínua} \emph{(combine
o~\cref{thm:b2:realfun:monotone} e o
\cref{cor:b2:realfun:nojumps})}.
\end{exercise}

\begin{exercise}[$\star\star$]\label{exo:b2:realfun:6}
(A média geométrica como limite) Para $x_i$ positivos e pesos
$\lambda_i$ somando $1$, prove que
\[
\lim_{p \to 0^+} \Bigl(\sum_i \lambda_i x_i^p\Bigr)^{1/p}
= \prod_i x_i^{\lambda_i} ,
\]
via $x_i^p = \eu^{p\ln x_i} = 1 + p\ln x_i + O(p^2)$, e deduza
a desigualdade MA--MG com pesos a partir do
\cref{ex:b2:realfun:powermeans}.
\end{exercise}

\begin{exercise}[$\star\star$]\label{exo:b2:realfun:7}
(Desigualdade da entropia) Usando a convexidade estrita de $t \mapsto t\ln t$,
prove que, para $p_i, q_i$ positivos com $\sum p_i = \sum q_i = 1$:
\[
\sum_i p_i \ln\frac{p_i}{q_i} \geq 0 ,
\]
com igualdade se e somente se $p = q$. \emph{(Escreva o membro da esquerda como $\sum q_i\,
\varphi\bigl(\frac{p_i}{q_i}\bigr)$ com $\varphi(t) = t\ln t$ e
aplique Jensen com pesos $q_i$.)}
\end{exercise}

\begin{exercise}[$\star\star\star$]\label{exo:b2:realfun:8}
(Convexidade no ponto médio) $f \colon I \to \R$ é \emph{convexa no ponto médio}
quando $f\bigl(\frac{x+y}{2}\bigr) \leq \frac{f(x) + f(y)}{2}$
sempre. Prove que uma função \emph{\omterm{def:b2:metric:continuity}{contínua}} convexa no ponto médio é
convexa. \emph{(Estabeleça a desigualdade de convexidade para pesos
diádicos $\frac{k}{2^m}$ por indução sobre $m$ e depois passe ao limite
usando densidade e \omterm{def:b2:metric:continuity}{continuidade}.)}
\end{exercise}

\begin{exercise}[$\star\star\star$]\label{exo:b2:realfun:9}
Seja $f$ convexa em $\intco{0}{+\infty}$ com $f(0) \leq 0$. Prove
que $x \mapsto \frac{f(x)}{x}$ é crescente em
$\intoo{0}{+\infty}$, e deduza que, para $f$ convexa com $f(0) =
0$: $f(x + y) \geq f(x) + f(y)$ para $x, y \geq 0$
(superaditividade).
\end{exercise}

\begin{exercise}[$\star$]\label{exo:b2:realfun:10}
Seja $f$ convexa em $I$ e $g$ convexa \emph{crescente} num
intervalo contendo $f(I)$. Prove que $g \circ f$ é convexa e
mostre, por um contraexemplo, que a monotonicidade de $g$ não pode ser
dispensada.
\end{exercise}

\begin{exercise}[$\star\star$]\label{exo:b2:realfun:11}
(Hermite--Hadamard) Seja $f$ convexa e \omterm{def:b2:metric:continuity}{contínua} em
$\intcc{a}{b}$. Prove
\[
f\Bigl(\frac{a+b}{2}\Bigr) \;\leq\; \frac{1}{b -
a}\int_a^b f(t)\,\dd t \;\leq\; \frac{f(a) + f(b)}{2} .
\]
\emph{(À esquerda: integre uma reta de suporte no ponto médio. À direita:
majore $f$ pela corda.)}
\end{exercise}

\begin{exercise}[$\star\star\star$]\label{exo:b2:realfun:12}
Prove que uma função convexa num intervalo \emph{\omterm{def:b2:metric:topology}{aberto}} $I$ é
localmente \omterm{def:b2:metric:continuity}{lipschitziana}: para todo segmento $\intcc{a}{b} \subseteq I$
e toda margem $\delta > 0$ com $\intcc{a - \delta}{b + \delta}
\subseteq I$, a restrição de $f$ a $\intcc{a}{b}$ é
\omterm{def:b2:metric:continuity}{lipschitziana}, com constante
$\max\Bigl(\bigl|\frac{f(a) - f(a - \delta)}{\delta}\bigr|,
\bigl|\frac{f(b + \delta) - f(b)}{\delta}\bigr|\Bigr)$
\emph{(prenda toda inclinação de corda entre essas duas pela desigualdade
das inclinações)}.
\end{exercise}

\section{Problema: a caixa de ferramentas da convexidade}

Uma definição --- a corda acima do gráfico --- gera toda a
caixa de ferramentas das desigualdades clássicas. Este problema de fim de semana
a constrói em ordem lógica: critérios de convexidade e Jensen estrito,
depois Young, Hölder e Minkowski (as certidões de nascimento das
\omterm{def:b2:nvs:norm}{normas} $p$), a cadeia \omterm{def:b2:metric:complete}{completa} das médias de potência, do mínimo ao
máximo, e dois dividendos de coroa --- a desigualdade de Carleman e
Hölder lido como uma dualidade. Tudo é demonstrado; nada é
importado.

\begin{problem}[{Problema de fim de semana --- Young, Hölder, Minkowski
e a cadeia das médias de potência}]
\label{pb:b2:realfun:1}
Em todo o problema, $p, q > 1$ são \emph{expoentes conjugados}: $\frac1p +
\frac1q = 1$; os vetores são $a = (a_1, \dots, a_n) \in \R^n$;
os pesos $\lambda_i > 0$ satisfazem $\sum_i\lambda_i = 1$.

\medskip
\noindent\textbf{Parte I --- Critérios e Jensen estrito.}
\begin{enumerate}
  \item Seja $f$ derivável num intervalo $I$. Prove que
        $f$ é convexa se e somente se $f'$ é crescente
        \emph{(uma direção passando ao limite na
        desigualdade das inclinações \cref{lem:b2:realfun:slopes}; a
        outra pelo teorema do valor médio)}. Deduza o critério
        $C^2$: $f'' \geq 0$.
  \item Suponha $f''> 0$ em $I$. Prove que $f$ é
        \emph{estritamente} convexa (desigualdade estrita para $x \neq y$
        e $\lambda \in \intoo01$) e que uma função estritamente convexa
        satisfaz a desigualdade de Jensen
        (\cref{thm:b2:realfun:convexreg}~(3)) com igualdade
        \emph{somente} quando todos os $x_i$ coincidem.
  \item Certifique as matérias-primas da caixa: $-\ln$ é estritamente
        convexa em $\intoo{0}{+\infty}$; $t \mapsto t^r$ é
        estritamente convexa aí para $r > 1$ e estritamente côncava
        para $0 < r < 1$; $\exp$ é estritamente convexa em $\R$.
  \item (Desigualdade de Young) Para $a, b \geq 0$, prove
        \[
        ab \;\leq\; \frac{a^p}{p} + \frac{b^q}{q},
        \]
        com igualdade se e somente se $a^p = b^q$ \emph{(aplique
        a concavidade de $\ln$ aos dois pontos $a^p, b^q$
        com pesos $\frac1p, \frac1q$)}.
  \item Redemonstre MA--MG com pesos em uma linha a partir da
        concavidade de $\ln$:
        \[
        \prod_i x_i^{\lambda_i} \leq \sum_i\lambda_ix_i \qquad
        (x_i > 0),
        \]
        com o caso de igualdade; compare com o caminho por limites do
        \cref{exo:b2:realfun:6}.
\end{enumerate}

\medskip
\noindent\textbf{Parte II --- Hölder e Minkowski.} Escreva
$\norm{a}_p = \bigl(\sum_i \abs{a_i}^p\bigr)^{1/p}$ e
$\norm{a}_\infty = \max_i\abs{a_i}$.
\begin{enumerate}[resume]
  \item (Hölder) Prove
        \[
        \sum_{i=1}^{n}\abs{a_ib_i} \;\leq\;
        \norm a_p\,\norm b_q ,
        \]
        com igualdade se e somente se os vetores $(\abs{a_i}^p)$ e
        $(\abs{b_i}^q)$ são proporcionais \emph{(normalize
        $\norm a_p = \norm b_q = 1$ e aplique Young termo a termo)}.
  \item Identifique os casos particulares: $p = q = 2$
        (Cauchy--Schwarz) e o par extremo $(p, q) = (1,
        \infty)$: enuncie e prove $\sum\abs{a_ib_i} \leq
        \norm a_1\norm b_\infty$.
  \item (Minkowski) Para $p \geq 1$, prove
        \[
        \norm{a + b}_p \leq \norm a_p + \norm b_p
        \]
        \emph{(escreva $\abs{a_i + b_i}^p \leq \abs{a_i +
        b_i}^{p-1}(\abs{a_i} + \abs{b_i})$ e aplique Hölder
        a cada produto)}. Conclua: $\norm\cdot_p$ é uma \omterm{def:b2:nvs:norm}{norma}
        em $\R^n$ para todo $p \in \intco{1}{+\infty}$,
        completando o quadro do~\cref{ch:b2:nvs}.
  \item Versões integrais: para $f, g$ \omterm{def:b2:metric:continuity}{contínuas} em
        $\intcc{a}{b}$, enuncie e prove Hölder e Minkowski
        para $\norm f_p = \bigl(\int_a^b\abs
        f^p\bigr)^{1/p}$ \emph{(mesmas demonstrações, com a
        positividade estrita da integral para a discussão do caso de
        igualdade)}.
  \item Prove a monotonicidade $\norm a_q \leq \norm a_p$ para $1
        \leq p \leq q$, o limite $\norm a_p \to
        \norm a_\infty$ quando $p \to \infty$, e a comparação
        recíproca com a constante ótima:
        \[
        \norm a_p \leq n^{\frac1p - \frac1q}\,\norm a_q
        \]
        \emph{(Hölder contra o vetor constante)}. Identifique
        os vetores que realizam cada igualdade.
  \item (Interpolação) Para $1 \leq p < r < q$ e $\theta \in
        \intoo01$ com $\frac1r = \frac\theta p +
        \frac{1-\theta}q$, prove
        \[
        \norm a_r \leq \norm a_p^{\theta}\,
        \norm a_q^{1-\theta}
        \]
        \emph{(aplique Hölder com os expoentes $\frac{p}{\theta
        r}$ e $\frac{q}{(1-\theta)r}$ a $\abs{a_i}^{\theta
        r}\abs{a_i}^{(1-\theta)r}$)}.
\end{enumerate}

\medskip
\noindent\textbf{Parte III --- A cadeia das médias de potência, \omterm{def:b2:metric:complete}{completa}.}
Para $p \neq 0$ ponha $M_p = \bigl(\sum_i\lambda_i
x_i^p\bigr)^{1/p}$ ($x_i > 0$) e $M_0 = \prod_i
x_i^{\lambda_i}$.
\begin{enumerate}[resume]
  \item Prove que $p \mapsto M_p$ é crescente em todo o
        $\R^*$: trate $p < q < 0$ pela identidade dos inversos
        $M_{-p}(x) = M_p(1/x)^{-1}$ e faça a ponte por $0$
        mostrando que $M_p \leq M_0 \leq M_q$ para $p < 0 < q$
        \emph{(aplique a concavidade de $\ln$ a $x_i^q$ e a
        desigualdade invertida para expoentes negativos)}.
  \item Prove os limites $M_p \to \max_i x_i$ quando $p \to
        +\infty$ e $M_p \to \min_i x_i$ quando $p \to -\infty$.
  \item Escreva a cadeia $\min \leq \mathrm{HM} \leq
        \mathrm{GM} \leq \mathrm{AM} \leq \mathrm{QM} \leq
        \max$ para pesos iguais e prove a consequência
        clássica: para $a_1, \dots, a_n$ positivos,
        \[
        \Bigl(\sum_i a_i\Bigr)\Bigl(\sum_i\frac1{a_i}\Bigr)
        \geq n^2 .
        \]
  \item Relacione médias e \omterm{def:b2:nvs:norm}{normas}: para pesos iguais $\lambda_i =
        \frac1n$, $M_p(x) = n^{-1/p}\norm x_p$. Concilie as
        duas monotonicidades --- as médias \emph{crescem} com $p$
        enquanto as \omterm{def:b2:nvs:norm}{normas} \emph{decrescem} (questão 10) --- em uma
        frase sobre o fator $n^{-1/p}$.
  \item Determine os casos de igualdade ao longo de toda a cadeia da
        questão 14 (pesos positivos): a igualdade em qualquer ponto
        força todos os $x_i$ iguais --- a convexidade estrita compensa.
\end{enumerate}

\medskip
\noindent\textbf{Parte IV --- Dividendos.}
\begin{enumerate}[resume]
  \item (Young com um botão de ajuste) Para $a, b \geq 0$ e $\varepsilon >
        0$, prove
        \[
        ab \leq \varepsilon\,\frac{a^p}{p} +
        \varepsilon^{-q/p}\,\frac{b^q}{q},
        \]
        e o caso cavalo de batalha $ab \leq \varepsilon a^2 +
        \frac{b^2}{4\varepsilon}$: o truque de absorção usado
        em toda a análise.
  \item (Rumo a Carleman) Seja $c_k = \frac{(k+1)^k}{k^{k-1}}$.
        Prove a identidade telescópica $\prod_{k=1}^{n}c_k =
        (n+1)^n$ e deduza, pela MA--MG aplicada aos números
        $c_ka_k$,
        \[
        (a_1a_2\cdots a_n)^{1/n} \leq
        \frac{1}{n(n+1)}\sum_{k=1}^{n} c_k a_k
        \qquad (a_k > 0).
        \]
  \item (Desigualdade de Carleman) Some sobre $n$, troque a
        ordem de somação (famílias positivas \omterm{def:b2:series:summable}{somáveis},
        \cref{thm:b2:series:fubini}) e use $\sum_{n \geq
        k}\frac{1}{n(n+1)} = \frac1k$ e $c_k/k = \bigl(1 +
        \frac1k\bigr)^k < \eu$ para concluir: para toda
        série convergente $\sum a_k$ de termos positivos,
        \[
        \sum_{n=1}^{\infty}(a_1a_2\cdots a_n)^{1/n}
        \;\leq\; \eu\sum_{k=1}^{\infty}a_k .
        \]
  \item Para $f$ \omterm{def:b2:metric:continuity}{contínua} e positiva em $\intcc{0}{1}$,
        prove
        \[
        \Bigl(\int_0^1 f\Bigr)\Bigl(\int_0^1\frac1f\Bigr)
        \geq 1,
        \]
        com igualdade se e somente se $f$ é constante
        \emph{(Cauchy--Schwarz em $\sqrt f\cdot\frac1{\sqrt
        f}$)}.
  \item (Geometria das bolas) Usando o caso de igualdade de
        Minkowski, mostre que, para $1 < p < \infty$, a esfera
        unitária de $\norm\cdot_p$ não contém segmento (a \omterm{def:b2:nvs:norm}{norma}
        é estritamente convexa no sentido do
        \cref{pb:b2:nvs:1}), ao passo que, para $p = 1$ e $p =
        \infty$, ela contém: exiba as partes retas.
\end{enumerate}

\medskip
\noindent\textbf{Parte V --- Dualidade e síntese.}
\begin{enumerate}[resume]
  \item (Hölder como dualidade) Prove que, para todo $a \in
        \R^n$,
        \[
        \norm a_p = \max_{\norm b_q \leq 1}\ \sum_i a_ib_i ,
        \]
        exibindo explicitamente um $b$ maximizante. (A \omterm{def:b2:nvs:norm}{norma} $p$
        é a \omterm{def:b2:linalg:dual}{dual} da \omterm{def:b2:nvs:norm}{norma} $q$ --- o germe em dimensão finita
        da dualidade $L^p$.)
  \item (Momentos) Seja $X$ uma variável aleatória que assume um número finito
        de valores positivos $x_i$ com probabilidades
        $\lambda_i$. Reformule a questão 12 como: $r \mapsto
        \E[X^r]^{1/r}$ é crescente --- a desigualdade
        dos momentos (Lyapunov), a ser reutilizada no
        \cref{ch:b2:randomvar}.
  \item Resolva com ferramentas nomeadas, em duas linhas cada: (i) para
        $a, b, c$ positivos: $a^3 + b^3 + c^3 \geq
        \frac{(a+b+c)^3}{9}$; (ii) para $x_1, \dots,
        x_n$ positivos: $\bigl(\sum_i\sqrt{x_i}\bigr)^2 \leq
        n\sum_i x_i$.
  \item (Síntese) Trace a genealogia em cinco frases: da definição por
        cordas ao lema das inclinações; das inclinações às retas de suporte e a
        Jensen; da concavidade de $\ln$ a Young, a Hölder, a
        Minkowski e às \omterm{def:b2:nvs:norm}{normas} $p$; de Jensen à cadeia das médias
        de potência e aos momentos; de MA--MG a Carleman. Nomeie os cumes
        (Hölder--Minkowski; Carleman) e diga para onde a
        caixa de ferramentas se dirige: os espaços $L^p$ do volume do terceiro ano de
        graduação, cujos axiomas são exatamente as questões 6 e 8.
\end{enumerate}
\end{problem}
